
Training a 7B parameter language model costs millions of dollars in compute. Yet foundation model developers routinely train on datasets where 30-40\% of tokens are duplicates, paying full compute cost for data that teaches nothing new. At current cloud GPU rates, removing duplicates from a 1.4T token dataset could save approximately \$1.5-2M per training run. Despite this inefficiency, scaling law research treats all data tokens as equivalent. The Chinchilla scaling laws predict compute-optimal model configurations by optimizing the ratio of parameters N to dataset size D, but the formulation L(N, D, C) contains no quality term. One trillion deduplicated tokens is treated identically to one trillion tokens with 40\% redundancy.

This omission has practical consequences. Practitioners know data quality matters. GPT-3 used fuzzy deduplication heuristics, The Pile emphasized domain diversity, and LLaMA combined filtering with careful corpus balancing. Yet these decisions remain heuristic. Without an operational definition of data quality Q(D) integrated into scaling laws, practitioners face an impossible optimization problem: how much should we invest in curation versus simply scaling compute? A research lab with a fixed budget must choose between spending on deduplication infrastructure versus additional GPU-hours. Current scaling laws provide no principled answer because they assume Q(D) is constant.

The problem is that ``data quality'' has remained frustratingly subjective. While everyone agrees that deduplicated, diverse, low-perplexity data is ``better,'' we lack a unified measurement framework that: (1) defines quality quantitatively across multiple dimensions, (2) validates that these dimensions correlate with information content, and (3) provides reproducible metrics suitable for scaling law integration.

The core insight of our work is that data quality is measurable through information density---the information each token provides---and this density can be quantified via four non-redundant components: deduplication ratio (redundancy removal), domain diversity (concept coverage), perplexity score (signal-to-noise), and token efficiency (signal density). We hypothesized that these metrics would correlate with information-theoretic measures of data content, and that deduplication would emerge as the strongest predictor given its prominence in practitioner heuristics.

We validate this hypothesis empirically on 12 controlled C4 subsets (120GB total) with systematic quality variations across all four dimensions. Our key findings:

\begin{enumerate}
\item All four quality components correlate strongly (r > 0.5, p < 0.01) with information density measured via token entropy, n-gram redundancy, and semantic diversity.
\item Deduplication ratio is the strongest predictor (r = 0.72, p = 0.001), explaining 52\% of information density variance alone.
\item Composite Q(D) explains 61\% of density variance (\(R^2\) = 0.61).
\item Measurement is reproducible (coefficient of variation < 10\% across independent resamples).
\end{enumerate}

These results establish an operational definition of data quality suitable for scaling law research. We provide a validated, reproducible Q(D) metric that enables the research community to empirically investigate questions like: ``Does a 20\% quality improvement reduce compute requirements for the same performance target?'' Our contribution is a measurement framework---a tool that quantifies the waste---not a complete scaling law theory.
